\section*{Chapter \ref{ch:g2:balance-equilibrium} --- Balance and Equilibrium}

\begin{solution}{exo:g2:balance-equilibrium:1}
The Earth pulls the lamp down; the desk pushes it up, exactly as
hard. The two \omterm{def:g2:pushes-and-pulls:force}{forces} tie --- \omterm{def:g2:balance-equilibrium:equilibrium}{equilibrium} --- so the lamp stays put.
\end{solution}

\begin{solution}{exo:g2:balance-equilibrium:2}
When the two teams pull equally hard in opposite directions: the
pulls cancel. The tie breaks the moment one team pulls harder than
the other.
\end{solution}

\begin{solution}{exo:g2:balance-equilibrium:3}
It hangs level --- equal weights at equal distances. When one twin
scoots farther out, that twin's side goes down: distance counts too.
\end{solution}

\begin{solution}{exo:g2:balance-equilibrium:4}
Close to the pivot, with Lina far out on her side. The rule: a
heavier rider near the middle can tie with a lighter rider far away
--- weight and distance both count.
\end{solution}

\begin{solution}{exo:g2:balance-equilibrium:5}
The ruler's \omterm{def:g2:balance-equilibrium:point}{balance point} is at its middle; the hammer's is near the
heavy head. Check: rest each on one finger and find the spot where
it lies level.
\end{solution}

\begin{solution}{exo:g2:balance-equilibrium:6}
They meet close to the brush --- the heavy end. The \omterm{def:g2:balance-equilibrium:point}{balance point}
sits toward the heavy side of an object.
\end{solution}

\begin{solution}{exo:g2:balance-equilibrium:7}
The wide pyramid: it can lean a long way before its weight hangs out
past its broad base. The narrow tower topples at the smallest lean.
\end{solution}

\begin{solution}{exo:g2:balance-equilibrium:8}
The pole slows every wobble, giving the walker time to shift back
over the rope. On a curb you stretch your arms out for the same
reason: slower wobbles, easier catching.
\end{solution}

\begin{solution}{exo:g2:balance-equilibrium:9}
To stand on the right foot alone, your body must lean to the right,
so that your weight sits over that foot --- like the tower keeping
its weight over its base. The wall is exactly where you need to
lean: it forbids the move, your weight stays hanging over the lifted
foot's empty place, and you tip.
\end{solution}

\begin{solution}{exo:g2:balance-equilibrium:10}
The crane is a giant seesaw with its tower as the pivot. Rider one:
the load, far out on the long front arm. Rider two: the concrete
blocks --- much heavier, so they may sit close, on the short back
arm. Heavy-close ties with light-far, exactly the seesaw rule; the
back arm is short because its rider is the heavyweight.
\end{solution}
